\section{总结与延伸}

\subsection{核心总结表}

\begin{table}[H]
\centering
\begin{tabular}{p{0.26\textwidth} p{0.24\textwidth} p{0.25\textwidth}}
\toprule
主题 & 核心结论 & 工程提示 \\
\midrule
Q-Learning 核心 & 学习最优 Q 函数，策略由 $\arg\max_a Q(s,a)$ 得出 & 准备好 replay buffer、target network、regularization \\
DQN & replay buffer + target network + $\epsilon$-greedy 是可伸缩版 Q-Learning & 保持 buffer size、采样多样性与 target network 更新频率 \\
稳定化 & Double Q、防止 reward hack 以及 reward shaping & 监控 TD error、reward variance；设置 rollback guardrail \\
工程部署 & 评估矩阵 + checklist + telemetry & 用 sim rollout + human-in-the-loop replay 验证 checkpoint，线上 guardrail 监控 Q drift \\
\bottomrule
\end{tabular}
\caption{Lecture 6 的方法与工程对照}
\end{table}

\subsection{进一步阅读}
\subsection{进一步阅读}

\begin{itemize}
    \item Mnih et al., \textit{Playing Atari with Deep Reinforcement Learning} (DQN, 2013)
    \item Van Hasselt et al., \textit{Deep Reinforcement Learning with Double Q-learning} (Double DQN, 2016)
    \item Lillicrap et al., \textit{Continuous Control with Deep Reinforcement Learning} (DDPG, 2016)
    \item Schaul et al., \textit{Prioritized Experience Replay} (ICLR 2016)
    \item Jaderberg et al., \textit{Population Based Training of Neural Networks} (NeurIPS 2017) — for lifecycle automation
\end{itemize}

\subsection{本章小结}

Q-Learning 的价值在于从价值函数推导策略，结合稳定化、数据治理和工程化部署可以让这一经典算法在现代 RL 系统中可靠地运行。

\subsection{行动提醒}

\begin{itemize}
    \item 每轮训练后将 evidence matrix 与 order-of-operations 做成文档，方便下一次 rollback。
    \item 监控 `TD error drift`、`reward variance`、`episode success rate`，其中任一指标触发 guardrail后立即 snapshot。
    \item 安排 quarterly replay audit（human review 20 failure trajectories）。
\end{itemize}

\end{document}
